\kchapter{CI/CD and GitOps}{argo}

\begin{objectives}
\begin{itemize}
  \item the difference between Continuous Integration and Continuous Delivery;
  \item push-based vs.\ pull-based (GitOps) deployment;
  \item a CI pipeline that tests, builds, scans, and publishes images;
  \item GitOps with Argo CD: Applications, sync, self-heal, rollback;
  \item environment promotion and progressive delivery.
\end{itemize}
\end{objectives}

\section{From commit to cluster}

\begin{center}
\resizebox{\linewidth}{!}{%
\begin{tikzpicture}[
  st/.style={fill=#1, text=white, rounded corners=5pt, minimum width=2.25cm,
             minimum height=1.35cm, align=center, font=\scriptsize\bfseries},
  >={Stealth[length=2.6mm]}]
  \node[st=k8sgray]   (c) at (0,0)    {\Large\faCodeBranch\\commit};
  \node[st=k8sblue]   (t) at (2.7,0)  {\Large\faVial\\test};
  \node[st=k8sblue]   (b) at (5.4,0)  {\Large\faDocker\\build image};
  \node[st=k8sblue]   (s) at (8.1,0)  {\Large\faIcon{shield-alt}\\scan};
  \node[st=k8sdeep]   (p) at (10.8,0) {\Large\faUpload\\push to\\registry};
  \node[st=k8sdeep]   (g) at (13.5,0) {\Large\faIcon{git-alt}\\update tag\\in Git};
  \node[st=k8sorange] (a) at (16.2,0) {\kimg[0.6cm]{argo}\\Argo CD\\syncs};
  \node[st=k8snavy]   (k) at (18.9,0) {\kimg[0.6cm]{k8s-logo-white}\\cluster};
  \foreach \x/\y in {c/t,t/b,b/s,s/p,p/g,g/a,a/k} \draw[->, thick, k8snavy] (\x) -- (\y);
  \draw[decorate, decoration={brace, amplitude=6pt}, k8sblue, thick]
       (-1.1,0.95) -- (14.6,0.95) node[midway, above=7pt, font=\small\bfseries, text=k8sblue]{CI (e.g.\ GitHub Actions)};
  \draw[decorate, decoration={brace, amplitude=6pt, mirror}, k8sorange, thick]
       (15.1,-0.95) -- (20.0,-0.95) node[midway, below=7pt, font=\small\bfseries, text=k8sorange]{CD (GitOps)};
\end{tikzpicture}}
\end{center}

\begin{itemize}
  \item \kterm{Continuous Integration (CI)}: every change is automatically
        tested, built into an image, scanned, and published.
  \item \kterm{Continuous Delivery (CD)}: every published version can be deployed
        to an environment automatically and safely.
\end{itemize}

\section{Push vs.\ pull}

\begin{center}
\ktablestart
\begin{tabularx}{\linewidth}{P{2.4cm} L L}
\kheader \kth{} & \kth{Push (CI runs \kc{kubectl}/\kc{helm})} & \kth{Pull / GitOps (agent in cluster)} \\
Who deploys & The CI runner, with cluster credentials. & A controller inside the cluster. \\
Credentials & CI must hold a powerful kubeconfig. & Cluster only needs read access to Git. \\
Drift & Manual changes go unnoticed. & Detected and (optionally) reverted. \\
Audit / rollback & Re-run old pipeline. & \kc{git log} / \kc{git revert}. \\
Tools & GitHub Actions, GitLab CI, Jenkins. & Argo CD, Flux. \\
\end{tabularx}
\end{center}

\begin{keyidea}
\textbf{GitOps}: the desired state of every environment lives in Git. An agent
in the cluster continuously makes the cluster match Git. To deploy, you merge a
pull request. To roll back, you revert it. Git becomes your audit log.
\end{keyidea}

\section{A CI pipeline with GitHub Actions}

\begin{yamlcode}{.github/workflows/ci.yaml}
name: ci
on:
  push:
    branches: [main]
  pull_request:

jobs:
  build:
    runs-on: ubuntu-latest
    permissions:
      contents: read
      packages: write
    steps:
      - uses: actions/checkout@v4
      - uses: actions/setup-python@v5
        with: { python-version: "3.12" }
      - run: pip install -r requirements.txt -r requirements-dev.txt && pytest -q

      - uses: docker/setup-buildx-action@v3
      - uses: docker/login-action@v3
        with:
          registry: ghcr.io
          username: ${{ github.actor }}
          password: ${{ secrets.GITHUB_TOKEN }}
      - uses: docker/build-push-action@v6
        with:
          push: ${{ github.event_name == 'push' }}
          tags: ghcr.io/${{ github.repository }}:${{ github.sha }}
          cache-from: type=gha
          cache-to: type=gha,mode=max
\end{yamlcode}

The capstone (Chapter~\ref{ch:capstone-devops}) extends this with scanning,
manifest validation, and an automatic GitOps update.

\section{GitOps with Argo CD}

\kicon{argo}~\textbf{Argo CD} is a CNCF-graduated GitOps controller with a
web UI.

\begin{bashcode}{terminal}
kubectl create namespace argocd
kubectl apply -n argocd -f https://raw.githubusercontent.com/argoproj/argo-cd/stable/manifests/install.yaml

kubectl -n argocd get secret argocd-initial-admin-secret \
  -o jsonpath='{.data.password}' | base64 -d; echo
kubectl -n argocd port-forward svc/argocd-server 8443:443
# open https://localhost:8443  (user: admin)
\end{bashcode}

\subsection{The Application object}

\begin{yamlcode}{argocd-app.yaml}
apiVersion: argoproj.io/v1alpha1
kind: Application
metadata:
  name: shop-prod
  namespace: argocd
spec:
  project: default
  source:
    repoURL: https://github.com/example/shop-deploy.git
    targetRevision: main
    path: k8s/overlays/prod          # Kustomize, Helm chart, or plain YAML
  destination:
    server: https://kubernetes.default.svc
    namespace: shop-prod
  syncPolicy:
    automated:
      prune: true                    # delete objects removed from Git
      selfHeal: true                 # revert manual changes in the cluster
    syncOptions:
      - CreateNamespace=true
    retry:
      limit: 5
      backoff: { duration: 10s, factor: 2, maxDuration: 3m }
\end{yamlcode}

\begin{center}
\ktablestart
\begin{tabularx}{\linewidth}{P{3.0cm} L}
\kheader \kth{Status} & \kth{Meaning} \\
\textcolor{k8sgreen}{\textbf{Synced}} & The cluster matches Git. \\
\textcolor{k8sorange}{\textbf{OutOfSync}} & Git and the cluster differ; Argo CD will sync (if automated) or wait for you. \\
\textcolor{k8sgreen}{\textbf{Healthy}} & Resources are running correctly (Deployments available, etc.). \\
\textcolor{k8sred}{\textbf{Degraded}} & Something is failing (CrashLoop, failed rollout). \\
\end{tabularx}
\end{center}

\begin{protip}
Use the \textbf{app-of-apps} pattern or an \textbf{ApplicationSet} to manage
many Applications (every service \(\times\) every environment) from a single
template, and keep your \emph{application source} repo separate from your
\emph{deployment config} repo or folder.
\end{protip}

\section{Environment promotion}

\begin{itemize}
  \item \textbf{Dev}: CI updates the image tag in \kc{overlays/dev} on every
        merge to \kc{main}; Argo CD auto-syncs.
  \item \textbf{Staging}: a pull request copies the tag to \kc{overlays/staging};
        automated tests run against it.
  \item \textbf{Production}: a reviewed, approved pull request promotes the same
        immutable image (same digest) to \kc{overlays/prod}.
\end{itemize}

Tools such as \textbf{Argo CD Image Updater}, \textbf{Kargo}, or \textbf{Flux}
image automation can create these promotion commits for you.

\section{Progressive delivery}

\textbf{Argo Rollouts} replaces a Deployment with a \kc{Rollout} that does
canary or blue/green releases, shifting traffic in steps and checking metrics
(e.g.\ Prometheus error rate) before each step. If metrics degrade, it aborts
and rolls back automatically.

\begin{yamlcode}{rollout-canary.yaml}
apiVersion: argoproj.io/v1alpha1
kind: Rollout
metadata:
  name: api
spec:
  replicas: 5
  strategy:
    canary:
      steps:
        - setWeight: 10
        - pause: { duration: 5m }
        - analysis:
            templates:
              - templateName: error-rate-below-1pct
        - setWeight: 50
        - pause: { duration: 10m }
  selector:
    matchLabels: { app: api }
  template:
    metadata:
      labels: { app: api }
    spec:
      containers:
        - name: api
          image: ghcr.io/example/api:1.5.0
\end{yamlcode}

\begin{recap}
\begin{itemize}
  \item CI: test, build, scan, publish. CD: deploy safely and repeatably.
  \item GitOps keeps environments in Git and lets an in-cluster agent pull changes.
  \item Argo CD Applications sync, prune, and self-heal.
  \item Promote the \emph{same} image through environments via pull requests.
\end{itemize}
\end{recap}
